\documentclass[a4paper]{article}
\usepackage{geometry}
\usepackage{graphicx}
\usepackage{natbib}
\usepackage{amsmath}
\usepackage{amssymb}
\usepackage{amsthm}
\usepackage{paralist}
\usepackage{epstopdf}
\usepackage{tabularx}
\usepackage{longtable}
\usepackage{multirow}
\usepackage{multicol}
\usepackage[hidelinks]{hyperref}
\usepackage{fancyvrb}
\usepackage{algorithm}
\usepackage{algorithmic}
\usepackage{float}
\usepackage{paralist}
\usepackage[svgname]{xcolor}
\usepackage{enumerate}
\usepackage{array}
\usepackage{times}
\usepackage{url}
\usepackage{fancyhdr}
\usepackage{comment}
\usepackage{environ}
\usepackage{times}
\usepackage{textcomp}
\usepackage{caption}
\usepackage{color}
\usepackage{xcolor}



\urlstyle{rm}

\setlength\parindent{0pt} % Removes all indentation from paragraphs
\theoremstyle{definition}
\newtheorem{definition}{Definition}[]
\newtheorem{conjecture}{Conjecture}[]
\newtheorem{example}{Example}[]
\newtheorem{theorem}{Theorem}[]
\newtheorem{lemma}{Lemma}
\newtheorem{proposition}{Proposition}
\newtheorem{corollary}{Corollary}

\floatname{algorithm}{Procedure}
\renewcommand{\algorithmicrequire}{\textbf{Input:}}
\renewcommand{\algorithmicensure}{\textbf{Output:}}
\newcommand{\abs}[1]{\lvert#1\rvert}
\newcommand{\norm}[1]{\lVert#1\rVert}
\newcommand{\RR}{\mathbb{R}}
\newcommand{\CC}{\mathbb{C}}
\newcommand{\Nat}{\mathbb{N}}
\newcommand{\br}[1]{\{#1\}}
\DeclareMathOperator*{\argmin}{arg\,min}
\DeclareMathOperator*{\argmax}{arg\,max}
\renewcommand{\qedsymbol}{$\blacksquare$}

\definecolor{dkgreen}{rgb}{0,0.6,0}
\definecolor{gray}{rgb}{0.5,0.5,0.5}
\definecolor{mauve}{rgb}{0.58,0,0.82}

\newcommand{\Var}{\mathrm{Var}}
\newcommand{\Cov}{\mathrm{Cov}}

\newcommand{\vc}[1]{\boldsymbol{#1}}
\newcommand{\xv}{\vc{x}}
\newcommand{\Sigmav}{\vc{\Sigma}}
\newcommand{\alphav}{\vc{\alpha}}
\newcommand{\muv}{\vc{\mu}}

\newcommand{\red}[1]{\textcolor{red}{#1}}

\def\x{\mathbf x}
\def\y{\mathbf y}
\def\w{\mathbf w}
\def\v{\mathbf v}
\def\E{\mathbb E}
\def\V{\mathbb V}

% TO SHOW SOLUTIONS, include following (else comment out):
\newenvironment{soln}{
    \leavevmode\color{blue}\ignorespaces
}{}

\newcommand \mle [1]{{\hat #1}^{\rm MLE}}



\hypersetup{
%    colorlinks,
    linkcolor={red!50!black},
    citecolor={blue!50!black},
    urlcolor={blue!80!black}
}

\geometry{
  top=1in,            % <-- you want to adjust this
  inner=1in,
  outer=1in,
  bottom=1in,
  headheight=3em,       % <-- and this
  headsep=2em,          % <-- and this
  footskip=3em,
}


\pagestyle{fancyplain}
\lhead{\fancyplain{}{Homework 2}}
\rhead{\fancyplain{}{CS 760 Machine Learning}}
\cfoot{\thepage}

\newcommand{\R}{\mathbb R}
\usepackage{enumitem}
\usepackage{bm}
\def\*#1{\mathbf{#1}}

\title{\textsc{Homework 3}} % Title

%%% NOTE:  Replace 'NAME HERE' etc., and delete any "\red{}" wrappers (so it won't show up as red)

\author{
\red{$>>$NAME HERE$<<$} \\
\red{$>>$ID HERE$<<$}\\
} 

\date{}

\begin{document}

\maketitle 


\textbf{Instructions:} 

Use this LaTeX file as a template to develop your homework and submit it as a single PDF file on Gradescope. 
For the programming component, you can use any programming language, although the teaching staff will be most able to assist with Python and specifically the {\tt sklearn} and {\tt torch} packages; 
you can use any functionality for Question~\ref{sec:training}.
Put all code used for the assignment at the end of the document in whatever way is easiest but still legible.

\section{Short proofs}

\subsection{Multi-class prediction with the 0-1 loss}

Suppose the world generates a single observation $x \sim \mbox{multinomial}(\theta)$, where the parameter vector $\theta=(\theta_1, \ldots, \theta_k)$ with $\theta_i\ge 0$ and $\sum_{i=1}^k \theta_i=1$.  Note $x \in \{1, \ldots, k\}$.
You know $\theta$ and want to predict $x$. 
Call your prediction $\hat x$.  What is your expected 0-1 loss: 
$$\E 1_{\hat x \neq x}$$
using the following two prediction strategies respectively?  Prove your answer.

\subsubsection{Strategy 1 [4 points]}
$\hat x \in \argmax_x \theta_x$, the outcome with the highest probability.

{\color{blue} Your answer here.}

\subsubsection{Strategy 2 [6 points]}

You mimic the world by generating a prediction $\hat x \sim \mbox{multinomial}(\theta)$.  (Hint: your randomness and the world's randomness are independent)

{\color{blue} Your answer here.}

\subsection{Predicting optimally under a weighted multi-class loss [8 points]}

Like in the previous question, 
the world generates a single observation $x \sim \mbox{multinomial}(\theta)$.
Let $c_{ij} \ge 0$ denote the loss you incur, if $x=i$ but you predict $\hat x=j$, for $i,j \in \{1, \ldots, k\}$.
$c_{ii}=0$ for all $i$.
This is a way to generalize different costs on false positives vs. false negatives from binary classification to multi-class classification.
You want to minimize your expected loss:
$$\E c_{x \hat x}.$$
Derive the optimal prediction $\hat x$.

{\color{blue} Your answer here.}

\subsection{Brute-forcing $k$-means [7 points]}

Consider the following simple brute-force algorithm for minimizing the $k$-means
objective: enumerate all the possible partitionings of the $n$ points into $k$
clusters. For each possible partitioning, compute the optimal centers $c_1,\dots,c_k$ by taking the mean of the points in each cluster and compute the
corresponding $k$-means objective value. Output the best clustering found. This
algorithm is guaranteed to output the optimal set of centers.
For the case $k = 2$, argue that the running time of this brute-force algorithm is exponential in the number of data points $n$.
	
{\color{blue} Your answer here.}

\newpage
\section{Convolutional neural networks}

In this question, you will apply convolutions to the following input:

\begin{tabular}{|c|c|c|c|}
	\hline
	3 & 7 & 2 & 5 \\
	\hline
	6 & 4 & 9 & 3 \\
	\hline
	1 & 8 & 5 & 6 \\
	\hline
	9 & 2 & 3 & 1 \\
	\hline
\end{tabular}

\subsection{A basic 2D convolution [5 points]}

What is the output when a convolution with the following $2\times 2$ filter is applied to the above input with stride 2, dilation 1, and no padding?

\begin{tabular}{|c|c|}
	\hline
	0.1 & 0.5 \\
	\hline
	0.2 & 0.3 \\
	\hline
\end{tabular}

{\color{blue} Your answer here.}

\subsection{Dilated convolutions [5 points]}\label{sec:dilated}

What is the output when the same filter is applied to the same input but with stride 1, dilation 2, and no padding?

{\color{blue} Your answer here.}

\subsection{Full matrix equivalent [10 points]}

Suppose the input is flattened into a 16 dimensional vector as follows:
\begin{tabular}{|c|c|c|c|c|c|c|}
	\hline
	3 & 7 & 2 & $\cdots$ & 2 & 3 & 1 \\
	\hline
\end{tabular}.
What is the matrix that would need to be multiplied by this vector such that when the product is reshaped~(by unflattening i.e. applying the inverse of the just-described flattening operation) into a $2\times2$ matrix, the output will be same as that of the dilated convolution from Question~\ref{sec:dilated}?

{\color{blue} Your answer here.}

\subsection{Convolution as regularization [5 points]}

Following the previous question, explain how a convolutional layer can be viewed as a regularization of a regular feed-forward layer.

{\color{blue} Your answer here.}

\newpage
\section{Neural network training}\label{sec:training}

Consider a two-layer feed-forward network with $h$ hidden units that takes $d$-dimensional inputs $\*x\in\R^d$ and outputs $\*{\hat y}=\textrm{softmax}(\*W_2\sigma(\*W_1\*x))$,
where $\*W_1\in\R^{h\times d}$ and $\*W_2\in\R^{k\times h}$ are weights and $\sigma:\R^h\mapsto[0,1]^h$ is the sigmoid activation function.
Suppose we want to train this network to minimize the negative log-likelihood objective $L(\*{\hat y},\*y)=-\*y^\top\log f(\*x)$ given one-hot-encoded labels $\*y\in\{0,1\}^k$.

\subsection{Gradient expressions [10 points]}
Derive the gradient of $L(\*{\hat y},\*y)$ w.r.t. $\*W_1$ and $\*W_2$, in terms of $\*W_1,\*W_2,\*x,\*y$, and $\*{\hat y}$.
You may use the following facts about matrix-vector derivatives without proof:
\begin{enumerate}
	\item $\frac{\partial L(\textrm{softmax}(\*z),\*y)}{\partial\*z}=\textrm{softmax}(\*z)-\*y$
	\item if $f:\R^m\mapsto\R$ and $\*u=\*W\*v$ for $\*W\in\R^{m\times n}$ and $\*v\in\R^n$ then $\frac{\partial f}{\partial\*v}=\*W^\top\frac{\partial f}{\partial\*u}$ and $\frac{\partial f}{\partial \*W}=\frac{\partial f}{\partial\*u}\*v^\top$
	\item if $f=g(\*u)$ and $\*u=\sigma(\*v)$ then $\frac{\partial f}{\partial\*v}=\frac{\partial f}{\partial\*u}\odot\sigma'(\*v)=\frac{\partial f}{\partial\*u}\odot\sigma(\*v)\odot(\*1-\sigma(\*v))$
\end{enumerate}

{\color{blue} Your answer here.}

\subsection{Backpropagation algorithm [10 points]}
Write down the quantities that need to be stored in the forward pass and the step-by-step (one matrix or vector operation per step) algorithm for computing the two gradients in the backward pass.

{\color{blue} Your answer here.}

\subsection{MNIST classifier [10 points]}\label{sec:mnist}

Train and test the above model with $h=128$ on flattened MNIST images using SGD with learning rate 1E-3 and batch size 1.
The data consists of 60K training and 10K test points of labeled images of the numbers 0-9, with pixel values normalized to be between 0 and 1 by dividing by 255 if needed.
It is easiest to obtain this data via 
\begin{enumerate}
	\item {\tt from torchvision import datasets, transforms}
	\item {\tt tt = transforms.ToTensor()}
	\item {\tt train = datasets.MNIST('./data', train=True, transform=tt, download=True)}
	\item {\tt test = datasets.MNIST('./data', train=False, transform=tt, download=True)}
\end{enumerate}
Initialize the weights at each layer by setting the entries by sampling i.i.d. from the uniform distribution on $[\pm\tfrac1{\sqrt k}]$, where $k$ is the input dimension of that layer.
Train for five epochs and report (1)~the average loss observed across each epoch, (2)~the average training accuracy observed across each epoch, and (3)~the test accuracy after training.

{\color{blue} Your answer here.}

\subsection{Normalization [5 points]}

Normalize the data by subtracting the mean pixel value and dividing by the standard deviation of the pixel values.
Report (1)~the mean, (2)~the standard deviation, and (3)~everything reported from Question~\ref{sec:mnist}.
Briefly comment on the effect of doing normalization.

{\color{blue} Your answer here.}

\subsection{Activation functions [5 points]}\label{sec:activation}

Replace the sigmoid activation by a ReLU and report everything reported in Question~\ref{sec:mnist}.
Briefly comment on the effect of doing switching to a different activation function.

{\color{blue} Your answer here.}

\subsection{Confusion matrix [5 points]}\label{sec:confusion}

Plot the confusion matrix of the model trained in Question~\ref{sec:activation}, making sure to label the individual matrix entries, and report (1)~the most mislabeled ground truth class and (2)~the class pair causing the most model confusion and which class is being confused for which other one in that case.

{\color{blue} Your answer here.}
	
\subsection{Data augmentation [5 points]}

Would it be reasonable to do data augmentation via flipping the input image along the vertical axis or rotating the image 180 degrees?
Explain why or why not.

{\color{blue} Your answer here.}


\end{document}
